\documentclass[12pt]{article}
\usepackage{amsmath}
\usepackage{amssymb}
\usepackage{amsthm}
\usepackage{geometry}
\geometry{a4paper, margin=1in}

\theoremstyle{definition}
\newtheorem{definition}{Definition}
\newtheorem{concept}{Concept}

\theoremstyle{plain}
\newtheorem{theorem}{Theorem}
\newtheorem{lemma}[theorem]{Lemma}

\title{Jigsaw \\ \small{A file splitting \& reconstruction utility}}
\author{Stochastic Batman}
\date{April 2, 2026}

\begin{document}
	
	\maketitle
	
	\section{Motivation: The Threshold Problem}
	Consider a high-stakes scenario where an invaluable secret - such as a proprietary recipe, a cryptographic master key, or a highly confidential file - must be protected. If the secret is stored locally in plaintext, it is vulnerable to a single point of failure via theft. Encrypting the secret with a master key only displaces the problem: who holds the master key? If a single individual holds it, the system is vulnerable to their loss, betrayal, or absence. 
	
	If we attempt to solve this by fracturing the key into $n$ literal fragments and distributing them among $n$ participants, we create an overly fragile system. If even one participant is unavailable (e.g., on vacation or lost their fragment), the remaining $n-1$ participants are completely locked out. 
	
	What is fundamentally required is a \textit{threshold scheme}. We need a mathematical mechanism to distribute a secret into $n$ shares such that any subset of $k$ shares (where $k \le n$) is completely sufficient to reconstruct the original secret, while any subset of $k-1$ shares yields absolutely zero information about the secret.
	
	\section{Intuitive Explanation}
	To achieve this, we can exploit the geometric properties of polynomials. 
	
	Imagine a blank two-dimensional plane. If we place a single point on this plane, there are infinitely many straight lines that can pass through it. However, if we place \textit{two} distinct points, there is exactly one unique straight line that connects them. 
	
	We can extend this logic to higher degrees:
	\begin{itemize}
		\item 2 points uniquely determine a line (a polynomial of degree 1).
		\item 3 points uniquely determine a parabola (a polynomial of degree 2).
		\item In general, $k$ distinct points uniquely determine a polynomial of degree $k-1$.
	\end{itemize}
	
	Shamir's Secret Sharing leverages this exact property. If we want a threshold of $k$ participants to unlock the secret, we generate a random polynomial of degree $k-1$. We mathematically embed our secret as the $y$-intercept of this polynomial (the value of the polynomial when $x = 0$). 
	
	We then evaluate this polynomial at $n$ different non-zero $x$-coordinates, generating $n$ distinct points (shares). We distribute one point to each participant. 
	If $k$ participants pool their points together, they have enough geometric constraints to trace the exact curve back to the $y$-axis, thereby revealing the secret $y$-intercept. If a group only musters $k-1$ points, the curve remains underdetermined; there are infinitely many valid polynomials of degree $k-1$ passing through those $k-1$ points, and consequently, infinitely many possible $y$-intercepts. 
	
	To make this practically secure and suited for computers, we perform all of this geometry not over the infinite continuous space of real numbers (which suffers from floating-point rounding errors and leaks partial information like parity), but over a constrained, wrap-around mathematical universe known as a \textit{Finite Field}.
	
	\section{Rigorous Mathematics}
	
	We formalize the intuitive geometry using the algebraic structures of finite fields and polynomial rings.
	
	\subsection{Notation and Preliminary Concepts}
	Let $p$ be a prime number such that $p > n$ and $p > S$, where $S$ is the secret to be shared. 
	Let $\mathbb{F}_p$ denote the finite field of order $p$, representing the set of integers modulo $p$. 
	Let $\mathbb{F}_p[x]$ denote the ring of polynomials in one variable $x$ with coefficients drawn from $\mathbb{F}_p$. 
	
	\begin{concept}[Modular Arithmetic and Finite Fields]
		In $\mathbb{F}_p$, the standard operations of addition, subtraction, and multiplication are performed modulo $p$. Because $p$ is prime, every non-zero element $a \in \mathbb{F}_p$ has a unique multiplicative inverse $a^{-1} \in \mathbb{F}_p$ such that $a \cdot a^{-1} \equiv 1 \pmod p$. This guarantees that division by any non-zero element is well-defined, establishing $\mathbb{F}_p$ as a field.
	\end{concept}
	
	\subsection{Definitions}
	
	\begin{definition}[Shamir's $(k, n)$ Secret Sharing Scheme]
		Given a secret $S \in \mathbb{F}_p$, a threshold $k$, and a total number of shares $n$ (with $k \le n < p$), the scheme proceeds as follows:
		\begin{enumerate}
			\item \textbf{Polynomial Generation:} Choose $k-1$ coefficients $a_1, a_2, \dots, a_{k-1}$ uniformly at random from $\mathbb{F}_p$. Define the polynomial $P \in \mathbb{F}_p[x]$ of degree at most $k-1$ as:
			$$ P(x) = S + a_1 x + a_2 x^2 + \dots + a_{k-1} x^{k-1} $$
			Note that the secret is embedded as the constant term, such that $P(0) = S$.
			\item \textbf{Share Distribution:} Generate $n$ distinct shares by evaluating the polynomial at $n$ distinct, non-zero points in $\mathbb{F}_p$. For $i = 1, 2, \dots, n$, the $i$-th share is the tuple $(x_i, y_i)$ where $x_i \neq 0$ and $y_i = P(x_i) \pmod p$.
		\end{enumerate}
	\end{definition}
	
	\subsection{Theorems and Lemmas (Without Proofs)}
	
	To reconstruct the polynomial from $k$ shares, we rely on polynomial interpolation over the finite field.
	
	\begin{theorem}[Lagrange Interpolation over $\mathbb{F}_p$]
		Given a set of $k$ distinct points $D = \{(x_1, y_1), (x_2, y_2), \dots, (x_k, y_k)\}$ where $x_i, y_i \in \mathbb{F}_p$ and $x_i \neq x_j$ for all $i \neq j$, there exists a unique polynomial $P \in \mathbb{F}_p[x]$ of degree at most $k-1$ such that $P(x_i) = y_i$ for all $1 \le i \le k$. 
		
		This polynomial can be explicitly constructed using the Lagrange basis polynomials $\ell_j(x)$:
		$$ P(x) = \sum_{j=1}^{k} y_j \ell_j(x) $$
		where each basis polynomial is defined as:
		$$ \ell_j(x) = \prod_{\substack{1 \le m \le k \\ m \neq j}} \frac{x - x_m}{x_j - x_m} \pmod p $$
	\end{theorem}
	
	Since the secret is defined as $P(0)$, participants possessing $k$ shares can bypass reconstructing the full polynomial and directly compute the secret by evaluating the Lagrange interpolation at $x = 0$:
	$$ S = P(0) = \sum_{j=1}^{k} y_j \prod_{\substack{1 \le m \le k \\ m \neq j}} \frac{-x_m}{x_j - x_m} \pmod p $$
	
	Finally, we formalize the assertion that fewer than $k$ shares reveal absolutely zero mathematical information about the underlying secret $S$.
	
	\begin{lemma}[Information-Theoretic Security / Perfect Secrecy]
		Let an adversary possess an arbitrary subset of $k-1$ shares. For any arbitrary candidate secret $S' \in \mathbb{F}_p$, there exists exactly one polynomial $Q \in \mathbb{F}_p[x]$ of degree at most $k-1$ that interpolates the $k-1$ shares \textit{and} the point $(0, S')$. 
		
		Because the $k-1$ random coefficients $a_1, \dots, a_{k-1}$ were chosen uniformly and independently from $\mathbb{F}_p$, every polynomial passing through the $k-1$ shares is equally probable. Consequently, the conditional probability distribution of the secret $S$, given the $k-1$ shares, is strictly uniform over $\mathbb{F}_p$. The adversary gains no statistical advantage in guessing $S$.
	\end{lemma}
	
	
	\section{Implementation over $\mathrm{GF}(2^8)$}
	
	The theoretical treatment in Section~3 uses $\mathbb{F}_p$ for a prime $p$.
	In practice, Jigsaw operates on \emph{files}, which are sequences of bytes.
	A byte is an integer in $\{0, 1, \dots, 255\}$ - exactly 256 values.
	The natural field to use is therefore $\mathrm{GF}(2^8)$, the unique finite
	field of order $2^8 = 256$.  All of Shamir's guarantees carry over
	unchanged; only the arithmetic rules differ.
	
	\subsection{From Integers to Binary Polynomials}
	
	Every element of $\mathrm{GF}(2^8)$ is represented as a degree-7
	polynomial with coefficients in $\mathbb{F}_2 = \{0, 1\}$:
	$$
	a_7 x^7 + a_6 x^6 + \cdots + a_1 x + a_0, \qquad a_i \in \{0,1\}.
	$$
	The eight bits of a byte are exactly the eight coefficients $a_7, \ldots, a_0$.
	The byte value $v \in \{0,\dots,255\}$ and the polynomial it encodes are
	related by:
	$$
	v = \sum_{i=0}^{7} a_i \, 2^i.
	$$
	
	\textbf{Example.}  The byte \texttt{0xb2} $= 178 = 10110010_2$ encodes
	$x^7 + x^5 + x^4 + x$.
	
	\subsection{Addition in $\mathrm{GF}(2^8)$: Bitwise XOR}
	
	Two elements are added by adding their coefficient polynomials over
	$\mathbb{F}_2$.  Because $1 + 1 = 0$ in $\mathbb{F}_2$ (no carry), this is
	precisely a bitwise XOR of the two bytes:
	$$ a(x) + b(x) = a \oplus b $$
	Subtraction is identical to addition in characteristic 2, since $-1 = 1$
	in $\mathbb{F}_2$.
	
	\subsection{The Irreducible Polynomial $\mathbf{0x11b}$}
	
	Multiplying two degree-7 polynomials can yield a result of degree up to 14,
	which no longer fits in a byte.  To stay within $\mathrm{GF}(2^8)$ we must
	\emph{reduce} the result modulo a degree-8 polynomial that is
	\emph{irreducible} over $\mathbb{F}_2$ - one that has no polynomial factors
	of lower positive degree (the polynomial analogue of a prime number).
	
	Jigsaw uses the polynomial:
	$$ m(x) = x^8 + x^4 + x^3 + x + 1 $$
	In binary its coefficients are $1\,0001\,1011$, which is the 9-bit number
	$\texttt{0x11b} = 283$.  This is the same polynomial used by the
	\textsc{aes} standard.
	
	\subsection{Multiplication in $\mathrm{GF}(2^8)$}
	
	Given two bytes $a$ and $b$, their product in $\mathrm{GF}(2^8)$ is:
	$$
	a \cdot b = \bigl(a(x)\cdot b(x)\bigr) \bmod m(x),
	$$
	where the intermediate polynomial multiplication uses XOR instead of
	ordinary addition (coefficients are in $\mathbb{F}_2$), and the modular
	reduction is polynomial long division with XOR.
	
	\textbf{Hardware-level view.}
	A single step of the reduction can be described concretely.  Let $c$ be the
	current partial result stored in a 9-bit register.  If the degree-8 bit
	(bit 8) is set, XOR in \texttt{0x11b} to cancel it:
	\begin{equation}
		c' = \begin{cases} c \oplus \texttt{0x11b} & c[8] = 1, \\ c & \text{otherwise} \end{cases}
		\label{eq:reduce}
	\end{equation}
	Repeating this after each left-shift step implements the full
	``multiply-then-reduce'' procedure.
	
	\subsection{The Generator $\mathbf{0x03}$ and Discrete Logarithms}
	
	The non-zero elements of $\mathrm{GF}(2^8)$ form a multiplicative cyclic
	group of order 255:
	$$ \mathrm{GF}(2^8)^* = \bigl\{\, g^0,\, g^1,\, g^2,\, \dots,\, g^{254} \,\bigr\} $$
	for some \emph{primitive element} (generator) $g$.  A primitive element is
	one whose successive powers cycle through all 255 non-zero values before
	returning to 1.
	
	Jigsaw uses $g = \texttt{0x03}$, which represents the polynomial $x + 1$.
	One can verify computationally (and it is a known fact for \texttt{0x11b})
	that the order of $\texttt{0x03}$ is exactly 255; i.e.,
	$\texttt{0x03}^{255} = 1$ and $\texttt{0x03}^{d} \ne 1$ for any proper
	divisor $d$ of 255.  The divisors of 255 are 1, 3, 5, 15, 17, 51, 85, so
	this requires checking seven conditions, all of which hold.
	
	\textbf{Why $g = \texttt{0x03}$?}  It is the smallest byte value (after
	\texttt{0x01}) that is primitive with respect to \texttt{0x11b}, making it
	the conventional choice.  Not every non-zero element is primitive:
	$\texttt{0x02}$ has order 255 only for some irreducible polynomials;
	$\texttt{0x03}$ is primitive for \texttt{0x11b}.
	
	Because $g$ is a generator, every non-zero element $a$ can be written as
	$a = g^{\log_g a}$.  This defines the \emph{discrete logarithm}
	$\log_g : \mathrm{GF}(2^8)^* \to \mathbb{Z}_{255}$.
	
	\subsection{Precomputed Tables: Turning Multiplication into Addition}
	
	Using the discrete logarithm, multiplication of two non-zero elements
	reduces to addition of integers modulo 255:
	$$
	a \cdot b = g^{\,(\log_g a + \log_g b) \bmod 255}.
	$$
	For zero we always have $0 \cdot b = 0$.
	
	This motivates two lookup tables of size 256 (or 512 for the exponential
	table, as explained below):
	
	\begin{itemize}
		\item \textbf{Log table} $L$: \quad $L[a] = \log_g a$ for $a \in \{1,\dots,255\}$; $L[0]$ is a sentinel (unused in valid multiplications).
		\item \textbf{Exp table} $ E[i] = \begin{cases} g^i & i \in \{0,\dots,254\}, \\ g^{i - 255} & i \in \{255,\dots,511\} \end{cases} $
	\end{itemize}
	
	\textbf{Why is the exp table doubled to 512 entries?}
	When computing $L[a] + L[b]$, the sum can reach up to $254 + 254 = 508$.
	Rather than computing a modular reduction $(L[a] + L[b]) \bmod 255$ on
	every multiplication (a branch or integer division), we extend the table to 512 entries.
	The sum then indexes directly into $E$ without any modular arithmetic:
	$$ a \cdot b = E[\,L[a] + L[b]\,] \qquad (a, b \ne 0)$$
	replacing a polynomial reduction by three table lookups and one addition.
	
	\subsection{Table Construction Algorithm}
	
	The tables are built by iteratively multiplying a running value by the
	generator using the reduction rule in Equation \eqref{eq:reduce}:
	
	\begin{enumerate}
		\item Initialise $v := 1$ (representing $g^0 = 1$).
		\item For $i = 0, 1, \dots, 254$:
		\begin{enumerate}
			\item Set $E[i] := v$ and $L[v] := i$
			\item Compute $v := v \ll 1$
			\item If bit 8 of $v$ is set, set $v := v \oplus \texttt{0x11b}$
			(reduce modulo $m(x)$)
		\end{enumerate}
		\item For $i = 255, \dots, 511$: $E[i] := E[i-255]$
		\item Set $L[0] := 0$ (guardrail; $\log(0)$ is undefined)
	\end{enumerate}
	
	After construction, $L$ and $E$ satisfy $E[L[a]] = a$ for all $a \ne 0$,
	confirming they are mutual inverses over the non-zero elements.
	
	\subsection{Connecting Back to Shamir's Scheme}
	
	With $\mathrm{GF}(2^8)$ arithmetic in hand, the Shamir scheme from
	Section 3 applies verbatim, with $\mathrm{GF}(2^8)$ substituted for
	$\mathbb{F}_p$:
	
	\begin{itemize}
		\item The \textbf{secret} $S$ is a single byte $\in \{0,\dots,255\}$.
		\item The \textbf{polynomial} $P(x) = S + a_1 x + \cdots + a_{k-1}x^{k-1}$
		has coefficients in $\mathrm{GF}(2^8)$; addition and multiplication
		use XOR and the table-based field multiplication respectively.
		\item \textbf{Shares} are pairs $(x_i, P(x_i)) \in \mathrm{GF}(2^8)^2$.
		Since the field has 255 non-zero elements, $n \le 255$ shares are
		supported naturally.
		\item \textbf{Reconstruction} applies Lagrange interpolation over
		$\mathrm{GF}(2^8)$: all additions become XOR, all divisions use the
		table-based inverse ($a^{-1} = E[255 - L[a]]$).
	\end{itemize}
	
	A file of length $\ell$ bytes is split by applying this procedure
	independently to each byte, producing $n$ share files each of length $\ell$.
	The precomputed tables are computed once, saved to disk, and reused for
	every subsequent split or join operation.
	
\end{document}